\section{本讲小结}
微分方程把局部规律延伸为过程函数。正确建模比形式求积分更先行：变量、符号、系统边界与初值应一起确定。线性化必须保留约束与守恒量，阻尼要区分黏滞阻力和干摩擦，质量流入还会带来额外动量需求。所有解都应回代方程，并检查初值、极限及物理适用区间。
\section{练习}
\begin{question}
\begin{drill}[垂直方向的追逐]\label{c4d1}
目标 $A$ 从原点沿 $+y$ 方向以速率 $v$ 运动。追逐者 $B$ 从 $(a,0)$ 出发，以恒定速率 $kv$ 始终朝向 $A$，$k>1$。求相遇时间。若相遇点高度恰为 $a$，求 $k$。
\end{drill}
\fig{c4d1q}{瞄准方向不断改变的追逐}{\begin{tikzpicture}
\draw[axis](0,0)--(4.5,0)node[right]{$x$};\draw[axis](0,0)--(0,3.4)node[above]{$y$};
\fill(0,2.6)circle(2pt)node[left]{$A$};\fill(2.5,.8)circle(2pt)node[right]{$B$};\draw[guide](0,2.6)--(2.5,.8);
\draw[force,blue](0,2.6)--(0,3.4)node[right]{$v$};\draw[force,red](2.5,.8)--(1.65,1.412)node[above right]{$kv$};
\node[below]at(3.5,0){初始 $B=(a,0)$};
\end{tikzpicture}}\answerlink{c4d1}
\end{question}
\begin{question}
\begin{drill}[单摆的方程与能量]\label{c4d2}
摆长 $l$，摆球质量 $m$，从小角度 $\theta_0$ 静止释放。写出精确方程及小角解，求最低点速率的精确表达式和小角首项，说明两者如何一致。
\end{drill}
\fig{c4d2q}{摆角、切向方向与最低点}{\begin{tikzpicture}
\coordinate(O)at(0,2.5);\coordinate(P)at(1.1,.2);\draw[world](O)--(P);\draw[guide](O)--(0,0);\fill(P)circle(3pt)node[right]{$m$};
\draw[force](P)--++(0,-.8)node[right]{$mg$};\draw(0,1.8)arc(-90:-64.44:.7)node[midway,below]{$\theta$};
\draw[force,blue](P)--++(-.44,.92)node[midway,right]{$T$};
\node[left]at(0,0){最低点};
\end{tikzpicture}}\answerlink{c4d2}
\end{question}
\begin{question}
\begin{drill}[由位置决定的作用力]\label{c4d3}
质量 $m$ 的质点在光滑直线上从 $x=0$ 静止出发，受 $F(x)=F_0(1-x/L)$，$F_0,L>0$。求 $0<x\le L$ 内的速度，判断速度何处最大，并求到达 $L$ 的时间。
\end{drill}
\fig{c4d3q}{力逐渐减小，直到终点为零}{\begin{tikzpicture}
\draw[axis](0,0)--(4.6,0)node[right]{$x$};\draw[axis](0,0)--(0,2.5)node[above]{$F$};
\draw[blue,thick](0,2)--(4,0);\node[left]at(0,2){$F_0$};\node[below]at(4,0){$L$};
\end{tikzpicture}}\answerlink{c4d3}
\end{question}
\begin{question}
\begin{drill}[阻尼振子的能量耗散]\label{c4d4}
欠阻尼振子满足 $m\ddot x+c\dot x+kx=0$，且最终停在平衡位置。初始 $x=x_0,v=v_0$。不求位移显式解，计算 $\int_0^\infty v(t)^2\dd t$，并说明该积分何时收敛。
\end{drill}
\fig{c4d4q}{有阻尼的弹簧振子及相对平衡位置的位移}{\begin{tikzpicture}

\draw[world](-.2,-.1)--(-.2,1.3);\draw[world](-.2,0)--(5,0);\draw[world,decorate,decoration={zigzag,segment length=5pt,amplitude=3pt}](-.2,.6)--(2.4,.6)node[midway,above]{$k$};\draw[world,fill=blue!6](2.4,.2)rectangle(3.2,1);\node at(2.8,.6){$m$};\draw[axis](1.5,-.7)--(5,-.7)node[right]{$x$};\draw[guide](1.7,-.9)--(1.7,1.3);\node[below]at(1.7,-.7){$0$};\draw[dim](1.7,-1.3)--(2.8,-1.3)node[midway,below]{$x$};\draw[flow](3.2,.6)--(4.5,.6)node[right]{$v>0$};\draw[force](2.4,1.2)--(1.2,1.2)node[left]{$cv$};\node at(3.6,1.8){示意 $x>0,\ v>0$};
\end{tikzpicture}}\answerlink{c4d4}
\end{question}
